\section{Conclusion}
\label{sec:conclusion}

This article established a formally verified property theory for division
and modulo generated by recursive quotient--remainder normalization. Within
the stated domains, the following results show that the normalization is
canonical and satisfies the expected algebraic and periodic laws:

\begin{equation*}
\begin{aligned}
& \forall\, a, b \in \mathbb{N} : b \neq 0 \\
& a < b \implies a \operatorname{mod} b = a &&\text{[Trivial Case]} \\
& a < b \implies a \operatorname{div} b = 0 &&\text{[Trivial Case]}
\end{aligned}
\end{equation*}

\begin{equation*}
\begin{aligned}
\forall\, n \in \mathbb{N} : n &\neq 0 \\
n \operatorname{mod} n &= 0 &&\text{[Identity]} \\
n \operatorname{div} n &= 1 &&\text{[Identity]}
\end{aligned}
\end{equation*}

\begin{equation*}
\begin{aligned}
\forall\, n \in \mathbb{N} &: \\
n \operatorname{mod} 1 &= 0 &&\text{[Division by One]} \\
n \operatorname{div} 1 &= n &&\text{[Division by One]}
\end{aligned}
\end{equation*}

\begin{equation*}
\begin{aligned}
\forall\, a, b \in \mathbb{Z} &: a \geq 0,\; b > 0 \\
a \operatorname{mod} b &= a \mathbin{\%}
  b &&\text{[Native Modulo Compatibility]}
\end{aligned}
\end{equation*}

\begin{equation*}
\begin{aligned}
\forall a,b,q,r \in \mathbb{Z} &: b \neq 0,\; a = bq + r \\
\operatorname{mod}(a + b, b) &= \operatorname{mod}(a, b) &&\text{[Linear Shift]} \\
\operatorname{div}(a + b, b) &= \operatorname{div}(a, b) + 1 &&\text{[Linear Shift]} \\
\operatorname{mod}(a - b, b) &= \operatorname{mod}(a, b) &&\text{[Linear Shift]} \\
\operatorname{div}(a - b, b) &= \operatorname{div}(a, b) - 1 &&\text{[Linear Shift]}
\end{aligned}
\end{equation*}

\begin{equation*}
\begin{aligned}
\forall a,b,q,r,m \in \mathbb{Z} &: b \neq 0,\; a = bq + r \\
\operatorname{mod}(a + m \cdot b, b) &= \operatorname{mod}(a, b) &&\text{[Linear Shift by Multiplier]} \\
\operatorname{div}(a + m \cdot b, b) &= \operatorname{div}(a, b) + m &&\text{[Linear Shift by Multiplier]} \\
\operatorname{mod}(a - m \cdot b, b) &= \operatorname{mod}(a, b) &&\text{[Linear Shift by Multiplier]} \\
\operatorname{div}(a - m \cdot b, b) &= \operatorname{div}(a, b) - m &&\text{[Linear Shift by Multiplier]}
\end{aligned}
\end{equation*}

\begin{equation*}
\begin{aligned}
\forall\, a,b,q_x,r_x,q_y,r_y &\in \mathbb{N},\; b \neq 0,\;
  a = bq_x + r_x = bq_y + r_y \\
\operatorname{DivMod}(a,b,q_x,r_x).\operatorname{solve}
  &= \operatorname{DivMod}(a,b,q_y,r_y).\operatorname{solve}
  &&\text{[Unique Remainder]}
\end{aligned}
\end{equation*}

\begin{equation*}
\begin{aligned}
\forall\, a, b \in \mathbb{Z} &: b \neq 0 \\
a \operatorname{mod} b
  &= ( a \operatorname{mod} b ) \operatorname{mod} b
  &&\text{[Modulo Idempotence]}
\end{aligned}
\end{equation*}

\begin{equation*}
\begin{aligned}
\forall\, a, b, c \in \mathbb{Z} &: b \neq 0 \\
( a + c ) \operatorname{mod} b
  &= ( a \operatorname{mod} b + c \operatorname{mod} b )
     \operatorname{mod} b &&\text{[Distributivity, Addition]} \\
( a + c ) \operatorname{div} b
  &= a \operatorname{div} b + c \operatorname{div} b
     + ( a \operatorname{mod} b + c \operatorname{mod} b )
       \operatorname{div} b &&\text{[Distributivity, Addition]} \\
( a + c) \operatorname{mod} b
  &= (a \operatorname{mod} b) + (c \operatorname{mod} b)
     - b \cdot (((a \operatorname{mod} b) \\
  &\qquad + (c \operatorname{mod} b)) \operatorname{div} b)
  &&\text{[Distributivity, Addition]}
\end{aligned}
\end{equation*}

\begin{equation*}
\begin{aligned}
\forall\, a, b, c \in \mathbb{Z} &: b \neq 0 \\
( a - c ) \operatorname{mod} b
  &= ( a \operatorname{mod} b - c \operatorname{mod} b )
     \operatorname{mod} b &&\text{[Distributivity, Subtraction]} \\
( a - c ) \operatorname{div} b
  &= a \operatorname{div} b - c \operatorname{div} b
     + ( a \operatorname{mod} b - c \operatorname{mod} b )
       \operatorname{div} b &&\text{[Distributivity, Subtraction]} \\
( a - c ) \operatorname{mod} b
  &= (a \operatorname{mod} b) - (c \operatorname{mod} b)
     - b \cdot (((a \operatorname{mod} b) \\
  &\qquad - (c \operatorname{mod} b)) \operatorname{div} b)
  &&\text{[Distributivity, Subtraction]}
\end{aligned}
\end{equation*}

\begin{equation*}
\begin{aligned}
\forall\, a, b, c \in \mathbb{Z} &: b \neq 0 \\
a \operatorname{mod} b = 0
  &\implies ( a + c ) \operatorname{mod} b = c \operatorname{mod} b
  &&\text{[Divisible-Base Shift Invariance]}
\end{aligned}
\end{equation*}

\begin{equation*}
\begin{aligned}
b &> 0,\quad 0 < k < b \\
k \operatorname{mod} b + (b - k) \operatorname{mod} b
  &= b &&\text{[Symmetrical Modulo Pairs]}
\end{aligned}
\end{equation*}

\begin{equation*}
\begin{aligned}
\forall\, a, b \in \mathbb{N} : b \neq 0 \\
a \operatorname{mod} b = b - 1
  &\implies (a + 1) \operatorname{mod} b = 0 &&\text{[Unit-Step Increment]} \\
a \operatorname{mod} b \neq b - 1
  &\implies (a + 1) \operatorname{mod} b = (a \operatorname{mod} b) + 1
  &&\text{[Unit-Step Increment]} \\
a \operatorname{mod} b = b - 1
  &\implies (a + 1) \operatorname{div} b = (a \operatorname{div} b) + 1
  &&\text{[Unit-Step Increment]} \\
a \operatorname{mod} b \neq b - 1
  &\implies (a + 1) \operatorname{div} b = a \operatorname{div} b
  &&\text{[Unit-Step Increment]}
\end{aligned}
\end{equation*}

\begin{equation*}
\begin{aligned}
\forall\, n, p \in \mathbb{N} : p > 1 \\
\exists!\, k \in [0, p) &: \operatorname{mod}(n + k,\; p) = 0
  &&\text{[Exactly One Zero per Block]}
\end{aligned}
\end{equation*}

Those formally verified properties are collected in
\href{https://github.com/thiagomata/prime-numbers/blob/master/src/main/scala/v1/chapter2/div/properties/Summary.scala}{\texttt{Summary.scala}}
and supported by the individual proof modules linked above. The recursive
formulation makes the proof structure transparent: normalize $(q,r)$ without
changing $a=bq+r$, extract quotient and remainder from the final state, then
derive the algebraic laws from that normal form.

Taken together, the canonical solution theorem, shift laws, arithmetic
identities, unit-step transition, and one-zero-per-period result
characterize both the normalized quotient--remainder state and its periodic
behavior on consecutive integers.

\section{Future Work}
\label{sec:future-work}

The recursive normalization argument used throughout this article
generalizes beyond the integers: the same shift-and-check invariant applies
to any Euclidean domain equipped with a well-founded remainder measure,
suggesting a more general recursive division theorem. The zero-density
results of Section~6.14 also invite a natural multiplicative extension: when
a block of consecutive integers is filtered by several pairwise coprime
moduli, the resulting density is expected to be the product of the
individual single-modulus densities, in the spirit of the Chinese Remainder
Theorem~[1]. Formalizing that multiplicative extension, and connecting the
recursive $\operatorname{DivMod}$ state to congruence-class arithmetic more
broadly, are natural next steps building on the identities established here.
